\subsection{Intro}

Recommended Reading: Probability and Random Processes (Grimmett/Stirzker)

Mathematical study of randomness (random $\leftrightarrow$ deterministic)

\subsubsection*{How do we model Randomness?}

Coin: $P(H)=P(T)=0.5$\\
Biased Coin: $P(H)=0.7, P(T)=0.3$\\
Dice: $P(1)=P(2)=\cdots=P(6)=\frac{1}{6}$\\
$P(1 \text{ or } 3)=P(\{1,3\})=P(1)+P(3)=\frac{1}{3}$\\
Power Law: $P(n)=\frac{6}{\pi^2}\frac{1}{n^2}$\\
$P(n\t{ is even})=P(\{2,4,6,\dots\})=\frac{1}{2}$\\
$P(\t{4 divides }n)=P(\{4,8,\dots\})$\\

\subsubsection*{Sample space and probability measure}

\begin{definition}[Sample Space]
    Let $\Omega$ be a set. Let $F$ be a collection of subsets of $\Omega$, which is a $\sigma$-algebra
    \begin{enumerate}
        \item $\emptyset \in F$
        \item If $A \in F, then A^c \in F$
        \item If $A_1,A_2,\dots \in F$ then $\bigcup\limits_n A_n \in F$
    \end{enumerate}
\end{definition}

\begin{definition}[Probability Measure]
    A Probability meassure on ($\Omega, F$) is a function\\
    $\mathbb{P}: F \mapsto [0,1]$ s.t.
    \begin{itemize}
        \item $\P(\emptyset) = 0,\;\;\P(\Omega)=1$
        \item If $A_1,A_2,\dots\in F$ is disjoint collection then $\P(\bigcup\limits_n A_n)=\sum\limits_n \P(A_n)$
    \end{itemize}
\end{definition}

\begin{exercise}
\phantom{text}
\begin{itemize}
    \item $A,B \in F$
    \begin{proof}
        $A\cap B = (A^c \cup B^c)^c$\\
        $A^c, B^c \in F \implies A^c \cup B^c \in F \implies (A^c \cup B^c)^c \in F$
    \end{proof}
    \item If $A,B \in F, A \subset B \implies \P(A) \leq \P(B)$\\
    $B=A\cup(B \setminus A) \implies \P(B)=\P(A)+\P(B\setminus A) \implies \P(A)\leq\P(B)$
\end{itemize}
\end{exercise}